\subsection{Zero-Gram off-diagonal vanishing and retained trace costs}
\label{sec:flag-retained-zero-gram-traces}

We continue with the original period family, including its
concrete standing hypotheses and absolutely convergent orbital
expansion. We allow distinct original words, with last-vector norms
\(-T,-T'\), and use the original compact types. No identification
of the Poincare maps with an additional arithmetic theta chain
is assumed in this subsection.

\begin{lemma}[Unequal-norm zero-Gram returns]
\label{lem:zero-gram-unequal-norm}
For a relevant return with \(G=0\), the two full flags have adapted
bases with common tail \(f_b=c e_b\), \(j\le b\le r\),
where \(j=r-\dim(X\cap Y)\) and \(c\bar c=T'/T\).
The leading \(j\)-by-\(j\) cross pairing is nondegenerate.
There is a common compact \(U(k)\) at each real place, fixing both
full flags pointwise and mapping isomorphically to each of their
positive rank-\(k\) quotient groups.
\end{lemma}
\begin{proof}
The common-tail argument of
Lemma~\ref{lem:rational-flag-common-tail} works verbatim up to
the norm equation: the skew-Hermitian successor relation and
the proof that neither negative last vector belongs to \(X+Y\)
use no equality of their norms. If \(X\cap Y\ne0\), the common
tail and \(z'=cz\) after adaptation follow as there, and
comparison of norms now gives \(c\bar c=T'/T\).
The same perpendicularity argument gives nondegeneracy of the
leading cross pairing.

If \(X\cap Y=0\), the argument excluding \(D_r=0\) uses only
the implication \(D_r=0\Rightarrow D_{r+1}=0\), which still
holds for the equation
\(|D_{r+1}|^2=TT'|D_r|^2\). The last vectors may then be made
perpendicular to the nondegenerate space \(X+Y\) by adapted
changes. Its complement has negative index one. The zero-Gram
equation becomes \(|B(z,z')|^2=TT'\). The orthogonal component of
\(z'\) to the negative line \(Ez\) has squared norm zero, so
vanishes by positivity. This gives the same norm equation
for the last-vector ratio, also when \(r=0\).

Split off the nondegenerate leading paired space of signature
\((j,j)\), the common negative line, and the remaining shared
isotropic space of dimension \(r-j\) together with isotropic
duals, chosen perpendicular to the preceding summands by
Hermitian Witt decomposition. Their complement \(W\) is positive
of rank \(k\). The group acting on \(W\) and fixing all the
other summands pointwise fixes both adapted frames, hence both
original frames. Projection from \(W\) to each positive
quotient \((X^\perp\cap z^\perp)/X\) and
\((Y^\perp\cap (z')^\perp)/Y\) is an isometric isomorphism.
This gives the two asserted identifications of compact groups.
\end{proof}

\begin{theorem}[Off-diagonal zero-Gram vanishing]
\label{thm:flag-offdiagonal-zero-gram}
If the two prescribed irreducible compact types differ at at
least one real place, every original zero-Gram return has
zero orbital operator. In particular the distinct-word
condition in the original DataNW gives the off-diagonal
vanishing on that concrete standing family.
\end{theorem}
\begin{proof}
An irrelevant return contributes zero by integration of the
nontrivial relative character over its compact arithmetic
unipotent quotient, as in the original relevance vanishing.
For a relevant return use the compact group \(C\) at the
place where the types differ from the preceding lemma.
It normalizes the intersection \(R_\infty\) of the two
real flag unipotent groups because it fixes both frames.
The relative character on \(R_\infty\) is trivial by
relevance. The original orbital integral over the arithmetic
stabilizer quotient is thus its finite covolume times
the integral over \(R_\infty\backslash G_\infty\), with
compatible measures. This is just Weil integration for
the original stabilizer; the integrand, not each seed,
descends. Absolute convergence in the original standing
regime justifies this operation.

Left multiplication by \(C\) acts on that quotient.
Its Haar-measure multiplier is one: conjugation on the
unipotent stabilizer has positive determinant modulus,
a continuous multiplicative character of a compact group.
On the first seed the action of \(C\) is the original
compact representation through the positive quotient.
Indeed decompose an element of \(C\) into the fixed Levi
lift of its action on that quotient and an element of
the flag unipotent. The latter fixes the entire first
frame, so its simple-root character is trivial.
For the second seed use the conjugate by the return and
the analogous decomposition. The two quotient maps are
induced by the complex-linear isometries of \(W\) in the
preceding lemma, so introduce no outer or conjugate
automorphism of the compact representation.

Changing variables by this compact action in the
quotient orbital integral gives the intertwining identity
for the two compact types on this one factor. By Schur's
lemma that operator is zero. Tensoring with the other
place factors, or applying Schur's lemma to the two
representations restricted to this factor, gives zero for
the global operator, with the original finite phase and
covolume unchanged. Distinctness in DataNW is the
inequality of the ordered type parameters at each real
place, hence the two irreducible representations are
inequivalent. This proves the assertion.
\end{proof}

\paragraph{Trace comparison on the actual retained columns.}
Write \(A_i\) for the normalized period maps of
Theorem~\ref{global-gram-entry-unfold}, using the antilinear
isometric column dictionary there. The index \(i\) can range
over any finite subfamily of the original words.
Let \(d_i\) denote their actual full column dimensions,
and \(J_i\) the isometric inclusions of the retained
source subspaces of ranks \(r_i\). Let \(\Pi\) be the
actual orthogonal target projection of rank \(D_0>0\).
Write \(R_0=\sum_i r_i>0\), and let \(s_i\ge0\) be
prescribed scalar weights on the period maps (take
\(s_i=1\) for the unweighted case). Set
\(F_i=s_i A_iJ_i\), \(F=(F_i)_i\), and
\[
 K=F^*F,\qquad L=F^*(1-\Pi)F\succeq0,\qquad
 K_0=K-L .
\]
The source Gram for the retained maps is \(\operatorname{diag}K_0\)
(block-diagonal compression, not the entrywise diagonal),
and the target Gram is \(\Pi FF^*\Pi\) on its retained range.
For any prescribed \(c>0\), define
\[
 {\cal D}_c(B)=
 \operatorname{tr}\big((c/R_0)I-\operatorname{diag}B\big)_+
 +\operatorname{tr}\big(B-(c/D_0)I\big)_+
\]
for a Hermitian operator \(B\) on the retained source sum.
The nonzero eigenvalues of the two Grams \(K_0\) and
\(\Pi FF^*\Pi\) coincide, with multiplicities; the remaining
zero eigenvalues do not contribute to the upper positive
trace at the strictly positive threshold. Hence the
normalized actual pair of positive traces is exactly
\({\cal D}_c(K_0)/c\).

\begin{lemma}[One trace-norm cost, including target leakage]
\label{lem:retained-gram-error-trace}
If \(K=B+E\) with \(B,E\) Hermitian, then
\[
 {\cal D}_c(K_0)\le {\cal D}_c(B)+\|E\|_1+\operatorname{tr}L.
 \tag{RT1}
\]
No positivity of the approximate Gram \(B\) is required.
\end{lemma}
\begin{proof}
The variational formula for the trace of the positive part
gives
\[
 \operatorname{tr}((c/R_0)I-\operatorname{diag}K_0)_+
 \le \operatorname{tr}((c/R_0)I-\operatorname{diag}B)_+
       +\operatorname{tr}E_-+\operatorname{tr}L.
\]
Indeed each block-diagonal contraction \(P\) testing the
first trace satisfies
\(-\operatorname{tr}(PE)\le\operatorname{tr}E_-\) and
\(\operatorname{tr}(PL)\le\operatorname{tr}L\).
For the target trace use \(K_0\le K\) and
\(\operatorname{tr}(K-(c/D_0)I)_+
\le\operatorname{tr}(B-(c/D_0)I)_++\operatorname{tr}E_+\).
The sum proves (RT1). This charges the signed error once,
rather than using two triangle inequalities with its full
trace norm.
\end{proof}

\begin{corollary}[The entire zero-Gram branch costs exponentially little]
\label{cor:retained-zero-gram-cost}
Use the concrete standing family of the diagonal zero-Gram
bound and Theorem~\ref{thm:flag-offdiagonal-zero-gram}.
Let \(B\) be the retained compression of the full weighted
period Gram with its nonidentity zero-Gram contributions
removed. Then
\[
 g_c+b_c\le
 \frac{{\cal D}_c(B)+\operatorname{tr}L}{c}
 +\frac{e^{-kf\log(k)/2}}c\sum_i s_i^2\sqrt{r_i d_i}.
 \tag{RT2}
\]
In particular for \(s_i=1\) and \(c=R_0\) the last term is
at most
\(e^{-kf\log(k)/2}\sqrt{\sum_i d_i/R_0}\).
\end{corollary}
\begin{proof}
The omitted operator is block diagonal by off-diagonal
vanishing. Its full diagonal block on the original column
space has Hilbert--Schmidt norm at most
\(\sqrt{d_i}\exp(-kf\log(k)/2)\) by the established
diagonal zero-Gram theorem. If desired take Hermitian
parts of both the remaining and omitted blocks; their sum
is Hermitian and this does not increase the norm bound.
The original column isometry preserves the bound.
Compression by \(J_i\) does not increase the
Hilbert--Schmidt norm. Its rank-\(r_i\) trace norm is
bounded by the square root of that rank times the
Hilbert--Schmidt norm. Summation over the diagonal blocks,
with the weights \(s_i^2\), bounds \(\|E\|_1\) in (RT1).
Cauchy--Schwarz gives the unweighted special case.
\end{proof}

This is a comparison on the actual period columns, for arbitrary
orthogonal source and target restrictions. In particular it requires
no cardinality estimate for all pairs of excluded columns.
It keeps the exact target leakage
\(\operatorname{tr}L=\sum_i s_i^2\|(1-\Pi)A_iJ_i\|_{\rm HS}^2\)
and the prescribed normalization \(c\).
No conclusion about the sign or size of the remaining nonzero-Gram
kernel is inferred; nor is the arithmetic metric and lattice
dictionary for the final theta maps inferred from an operator
compression. The Adriazola integral central-cluster projectors
can be used for the source cuts when the corresponding actual
arithmetic lattice admissibility has been established.
